\chapter{iOS Architecture}

\section{Application Structure}
The iOS app is implemented in:
\begin{itemize}[leftmargin=2em]
\item \texttt{\detokenize{naruto_app/naruto_app/}}
\end{itemize}

Top-level flow:
\begin{itemize}[leftmargin=2em]
\item \texttt{naruto\_appApp.swift} launches \texttt{ContentView.swift}.
\item Navigation routes users from Home to Mode Setup and then into game views.
\item \texttt{CameraGameView.swift} handles free, speed, and tutorial modes.
\item \texttt{BattleGameView.swift} combines camera inference with battle simulation.
\end{itemize}

\section{Core Runtime Managers}
\begin{table}[H]
\centering
\caption{Manager responsibilities in the runtime pipeline}
\label{tab:ios_managers}
\begin{tabular}{p{0.28\linewidth}p{0.62\linewidth}}
\toprule
Component & Responsibility \\
\midrule
CameraManager & Camera permission, session setup, frame streaming, orientation and mirroring management. \\
HandLandmarkDetector & Hand landmark extraction with MediaPipe Tasks and Vision fallback. \\
FaceDirectionEstimator & Mouth and face direction estimation for fireball and burning ash effects. \\
GestureRecognizer & Feature preprocessing, CoreML model inference, score formatting, top-k output. \\
JutsuManager & Hold logic, sequence tracking, mode-specific trigger decisions, effect state memory. \\
\bottomrule
\end{tabular}
\end{table}

\section{Model and Backend Fallback Strategy}
\begin{itemize}[leftmargin=2em]
\item Gesture model loading uses the unified 126-feature model exclusively; legacy one-hand and two-hand models were removed because their feature layouts are incompatible and a silent fallback would corrupt predictions. A missing or unloadable model now fails loudly (assertion + fault log).
\item The recognizer records a rolling average of end-to-end pipeline latency (landmarks $\rightarrow$ features $\rightarrow$ CoreML) and logs it every 120 frames via \texttt{os.log}, so real-time claims are measured.
\item Landmark detection attempts MediaPipe runtime task models and falls back to Vision requests if unavailable.
\item Face direction estimation similarly uses MediaPipe face mesh first with Vision fallback.
\end{itemize}

\section{Frame-to-Action Pipeline}
\begin{enumerate}[leftmargin=2em]
\item Camera delivers a sample buffer.
\item Hand and face landmarks are extracted.
\item Landmark vectors are canonicalized for mirror and orientation consistency.
\item Classifier predicts sign label and confidence.
\item JutsuManager processes hold and sequence state.
\item Active jutsu state updates overlays, particles, sounds, and game counters.
\end{enumerate}

\section{Rendering Stack}
\begin{itemize}[leftmargin=2em]
\item \texttt{CameraView.swift} overlays skeletons and effects above AVCapture preview.
\item CAEmitterLayer handles fire, lightning, and ash particle systems.
\item SpriteKit scenes handle water dragon and kuchiyose summon animations.
\item Overlay UI components display target sign guidance and progress strips.
\end{itemize}

\placeholderfigure{ios_architecture_diagram}{figures/screenshots/ios_architecture_diagram.png}{iOS architecture diagram placeholder.}
\placeholderfigure{camera_overlay_runtime}{figures/screenshots/runtime_prediction.png}{Runtime camera overlay screenshot placeholder.}
